\Article{31}
Ist die Anzahl der Größen \(x,y,z\) etc. etwas größer, so wird die numerische Bestimmung der Werte \(A,B,C\) etc. aus den Gleichungen \(\xi=0,\allowbreak\eta=0.\quad\zeta=0\) etc. durch gewöhnliche Elimination recht mühsam. Deshalb haben wir in Theoria Motus Corporum Coelestium, Art.\ 182, einen besonderen Algorithmus angegeben und in der \textit{Disquisitione de elementis ellipticis Palladis} (Comm. recent. Soc. Gotting Vol.\ I) ausführlich erläutert, durch den diese Arbeit so einfach wird, wie es die Sache erlaubt. Die Funktion \(\Omega\) ist nämlich in eine Form wie diese zu bringen:
\[
\frac{u^0u^0}{\mathfrak{A}^0}+\frac{u'u'}{\mathfrak{B}'}+\frac{u''u''}{\mathfrak{C}''}+\frac{u'''u'''}{\mathfrak{D}'''}+\text{etc.}+M
\]
Hierbei sind die Divisoren \(\mathfrak{A}^0,\mathfrak{B}',\mathfrak{C}'',\mathfrak{D}'''\) etc. bestimmte Größen, \(u^0,u',u'',u'''\) etc. dagegen lineare Funktionen von \(x,y,z\) etc.; von diesen ist jedoch die zweite, \(u'\), frei von \(x\), die dritte, \(u''\), frei von \(x\) und \(y\), die vierte, \(u'''\), frei von \(x,y\) und \(z\) und so fort, sodass die letzte, \(u^{(\pi-1)}\), nur die letzte der Unbestimmten \(x,y,z\) etc. enthält. Schließlich sind die Koeffizienten, mit denen \(x,y,z\) etc. in \(u^0,u',u''\) etc. jeweils multipliziert werden, jeweils gleich \(\mathfrak{A}^0,\mathfrak{B}',\mathfrak{C}''\) etc. Ist dies ausgeführt, so setzen wir \(u^0=0,\allowbreak u'=0,\allowbreak u''=0,\allowbreak u'''=0\) etc.; daraus lassen sich die Werte der Unbekannten \(x,y,z\) etc. in umgekehrter Reihenfolge am bequemsten gewinnen. Es scheint nicht nötig, den Algorithmus selbst, durch den diese Umformung der Funktion \(\Omega\) ausgeführt wird, hier noch einmal zu wiederholen.

Die unbestimmte Elimination, mit deren Hilfe die Gewichte dieser Bestimmungen gefunden werden müssen, erfordert aber eine noch weit umfangreichere Rechnung. In der Ta\SourceLeaf{043}%
t lässt sich das Gewicht der Bestimmung der letzten Unbekannten (die allein in die letzte \(u^{(\pi-1)}\) eingeht) nach dem in Theoria Motus Corporum Coelestium Bewiesenen leicht als gleich dem letzten Glied in der Reihe der Divisoren \(\mathfrak{A}^0,\mathfrak{B}',\mathfrak{C}''\) etc. finden. Deshalb haben mehrere Rechner, um jene mühsame Elimination zu vermeiden, beim Fehlen anderer Hilfsmittel den Ausweg gewählt, den mehrfach erwähnten Algorithmus mit veränderter Reihenfolge der Größen \(x,y,z\) etc. zu wiederholen, sodass jede nacheinander die letzte Stelle einnahm. Wir hoffen daher, dass es den Geometern willkommen sein wird, wenn wir hier eine neue, aus einer tieferen Untersuchung des Gegenstands gewonnene Methode zur Berechnung der Gewichte der Bestimmungen darlegen, die nichts weiter zu wünschen übrig zu lassen scheint.

\Article{32}
Setzen wir also (I)
\[
\begin{array}{rclllll}
u^0&=&\mathfrak{A}^0x&+\mathfrak{B}^0y&+\mathfrak{C}^0z&+\text{etc.}&+\mathfrak{L}^0\\
u'&=&&\mathfrak{B}'y&+\mathfrak{C}'z&+\text{etc.}&+\mathfrak{L}'\\
u''&=&&&\mathfrak{C}''z&+\text{etc.}&+\mathfrak{L}''\\
\text{etc.}
\end{array}
\]
Daraus ergibt sich identisch
\[
\begin{aligned}
\tfrac12d\Omega&=\xi\,dx+\eta\,dy+\zeta\,dz+\text{etc.}\\[3pt]
&=\frac{u^0du^0}{\mathfrak{A}^0}+\frac{u'du'}{\mathfrak{B}'}+\frac{u''du''}{\mathfrak{C}''}+\text{etc.}\\[3pt]
&=u^0\left(dx+\frac{\mathfrak{B}^0}{\mathfrak{A}^0}dy+\frac{\mathfrak{C}^0}{\mathfrak{A}^0}dz+\text{etc.}\right)\\
&\qquad+u'\left(dy+\frac{\mathfrak{C}'}{\mathfrak{B}'}dz+\text{etc.}\right)+u''(dz+\text{etc.})+\text{etc.}
\end{aligned}
\]
woraus wir schließen (II)
\[
\begin{aligned}
\xi&=u^0\\
\eta&=\frac{\mathfrak{B}^0}{\mathfrak{A}^0}u^0+u'\\
\zeta&=\frac{\mathfrak{C}^0}{\mathfrak{A}^0}u^0+\frac{\mathfrak{C}'}{\mathfrak{B}'}u'+u''\\
&\text{etc.}
\end{aligned}
\]
Wir nehmen an, dass sich daraus die folgenden Formeln ableiten lassen (III)
\SourceLeaf{044}
\[
\begin{aligned}
u^0&=\xi\\
u'&=A'\xi+\eta\\
u''&=A''\xi+B''\eta+\zeta\\
&\text{etc.}
\end{aligned}
\]
Subtrahieren wir nun vom vollständigen Differential der Gleichung
\[
\Omega=\xi(x-A)+\eta(y-B)+\zeta(z-C)+\text{etc.}+M
\]
die Gleichung
\[
\tfrac12d\Omega=\xi\,dx+\eta\,dy+\zeta\,dz+\text{etc.}
\]
so folgt
\[
\tfrac12d\Omega=(x-A)d\xi+(y-B)d\eta+(z-C)d\zeta+\text{etc.}
\]
Dieser Ausdruck muss mit dem folgenden, aus III hervorgehenden identisch sein:
\[
\begin{aligned}
&\frac{u^0}{\mathfrak{A}^0}.d\xi+\frac{u'}{\mathfrak{B}'}(A'd\xi+d\eta)\\
&\qquad+\frac{u''}{\mathfrak{C}''}(A''d\xi+B''d\eta+d\zeta)+\text{etc.}
\end{aligned}
\]
Daraus schließen wir (IV)
\[
\begin{array}{rcllll}
x&=&\dfrac{u^0}{\mathfrak{A}^0}&+A'.\dfrac{u'}{\mathfrak{B}'}&+A''.\dfrac{u''}{\mathfrak{C}''}&+\text{etc.}+A\\[7pt]
y&=&&\dfrac{u'}{\mathfrak{B}'}&+B''.\dfrac{u''}{\mathfrak{C}''}&+\text{etc.}+B\\[7pt]
z&=&&&\dfrac{u''}{\mathfrak{C}''}&+\text{etc.}+C\\[4pt]
\text{etc.}
\end{array}
\]
Setzen wir in diesen Ausdrücken für \(u^0,u',u''\) etc. ihre aus III entnommenen Werte ein, so ist die unbestimmte Elimination vollendet. Insbesondere haben wir zur Bestimmung der Gewichte (V)
\[
\begin{aligned}
\relax[\alpha\alpha]&=\frac1{\mathfrak{A}^0}+\frac{A'A'}{\mathfrak{B}'}+\frac{A''A''}{\mathfrak{C}''}+\frac{A'''A'''}{\mathfrak{D}'''}+\text{etc.}\\[4pt]
[\beta\beta]&=\phantom{\frac1{\mathfrak{A}^0}+{}}\frac1{\mathfrak{B}'}+\frac{B''B''}{\mathfrak{C}''}+\frac{B'''B'''}{\mathfrak{D}'''}+\text{etc.}\\[4pt]
[\gamma\gamma]&=\phantom{\frac1{\mathfrak{A}^0}+\frac1{\mathfrak{B}'}+{}}\frac1{\mathfrak{C}''}+\frac{C'''C'''}{\mathfrak{D}'''}+\text{etc.}\\
&\text{etc.}
\end{aligned}
\]
Die Einfachheit dieser Formeln lässt nichts zu wünschen übrig. Auch für die übrigen
\SourceLeaf{045}%
Koeffizienten \([\alpha\beta],[\alpha\gamma],[\beta\gamma]\) etc. ergeben sich ebenso einfache Formeln; da sie jedoch seltener gebraucht werden, verzichten wir darauf, sie hier anzugeben.

\clearpage
